\section{Clustering Methods}\label{sec:bg:clustering}

Clustering partitions a finite set of points \(\{\vect{z}_i\}_{i=1}^{N} \subset \R^{d}\) into groups so that points inside a group are more similar to one another than to points in other groups.
In this thesis the points are pulse embeddings and the groups are emitters within a window or operating modes.
The clustering step runs only at inference, after the encoder has been trained (\cref{sec:method:arch}).
Classical algorithms differ in how similarity is defined, whether the number of clusters must be known in advance, and how they treat outliers.
Centroid methods such as \(k\)-means, and mixture models such as a \gls{gmm}, require a prescribed number of groups and typically assign every point.
When that number is unknown and some points should remain unassigned, density-based clustering is a better match.

\subsection*{\Gls{dbscan}}

\Gls{dbscan} recovers clusters from local density rather than from a fixed set of centres \parencite{ester1996density}.
Two hyperparameters define the density criterion: a neighbourhood radius \(\varepsilon > 0\) and a minimum neighbourhood size \(\mathrm{minPts} \in \N\).
A point is a core point if at least \(\mathrm{minPts}\) points (including itself) lie inside its Euclidean ball of radius \(\varepsilon\).
Points that are not cores but fall inside some core's ball are border points.
The remainder are labelled noise.
Clusters grow by connecting core points that fall in each other's \(\varepsilon\)-neighbourhoods and attaching their border points.
The number of clusters is then an output of the run, not an input.

This matches the inference setting of \cref{eq:method:cluster_decoding}: within each window the number of active emitters (and, separately, of modes) is not known in advance, and sparse outliers can be left unassigned instead of distorting the partition.
The main practical cost is sensitivity to \(\varepsilon\).
A single global radius assumes that all clusters share a comparable density scale.
The method chapter uses \gls{dbscan} with frozen \((\varepsilon,\mathrm{minPts})\) on each branch embedding.

Applied to raw high-dimensional features, density thresholds degrade because neighbourhoods become less informative as \(d\) grows.
Deep clustering therefore first learns a lower-dimensional embedding in which the desired partitions are easier to separate, then runs a classical operator such as \(k\)-means or \gls{dbscan} in that space \parencite{xie2016unsupervised}.
How a recovered partition is scored against a reference is defined in \cref{sec:bg:cluster_metrics}.
\Cref{sec:bg:dl} reviews the representation-learning side of that pipeline.
